\documentclass[11pt]{article}

% ─── Packages ────────────────────────────────────────────────────────────────
\usepackage{amsmath,amssymb,amsthm}
\usepackage[margin=1in]{geometry}
\usepackage{hyperref}
\usepackage{graphicx}
\usepackage{booktabs}
\usepackage{array}
\usepackage{multirow}
\usepackage{xcolor}
\usepackage{authblk}
\usepackage{setspace}
\usepackage{caption}
\usepackage{subcaption}
\usepackage{enumitem}
\usepackage{microtype}
\usepackage{siunitx}

\hypersetup{
  colorlinks=true,
  linkcolor=blue,
  citecolor=blue,
  urlcolor=blue
}

% ─── Theorem environments ────────────────────────────────────────────────────
\theoremstyle{plain}
\newtheorem{theorem}{Theorem}[section]
\newtheorem{lemma}[theorem]{Lemma}
\newtheorem{proposition}[theorem]{Proposition}
\newtheorem{corollary}[theorem]{Corollary}
\theoremstyle{definition}
\newtheorem{definition}{Definition}[section]
\theoremstyle{remark}
\newtheorem{remark}{Remark}[section]

% ─── Title block ─────────────────────────────────────────────────────────────
\title{\textbf{Within-Class Geometric Diversity as a Practical\\
Mechanism for Coordination Entropy Tuning:\\
LiFSI in 2-MeTHF/DME and FEME/2-MeTHF\\
as Intra-Ether Approaches to SCE* = 1.466}}

\author[1]{Eric Robert Lawson}

\affil[1]{Independent Researcher, OrganismCore Initiative\\
\texttt{ORCID: 0009-0002-0414-6544}\\
Repository: \url{https://github.com/Eric-Robert-Lawson/attractor-oncology}}

\date{April 4, 2026\\
\small Pre-registration record:
\url{https://github.com/Eric-Robert-Lawson/attractor-oncology}\\
\small (timestamped commit history constitutes the priority
and pre-registration record)\\
\small Part of the SCE Framework preprint series.
Preprint~1 (foundation and fixed point):
same repository.\\
\small The full geometric structure of this series,
including the relationship between all candidate formulations,
mechanisms, and derivations, is navigable from the repository
index.}

% ─── Document ────────────────────────────────────────────────────────────────
\begin{document}

\maketitle

% ─── Abstract ────────────────────────────────────────────────────────────────
\begin{abstract}
The SCE Framework establishes $\mathrm{SCE}^{*} = 1.466$ as the Shannon
Coordination Entropy of the Li$^+$ first solvation shell at which
combined room-temperature (RT) and low-temperature (LT) Coulombic
efficiency (CE) is simultaneously maximised in lithium metal batteries.
Prior papers in this series introduced two orthogonal mechanisms for
reaching $\mathrm{SCE}^{*}$: ESP matching via cross-class co-solvents
(DME/TMP, Preprint~2) and anion reception via selective Lewis acid
addition (DOL/TMB, Preprint~3).

The present paper introduces a third mechanism: \textbf{within-class
geometric diversity}, in which two co-solvents from the same chemical
family (the ether class) but with structurally distinct coordination
geometries are mixed to raise SCE from below.  Cyclic ethers (2-MeTHF)
and linear ethers (DME, FEME) present geometrically distinguishable
Li$^+$ coordination modes despite belonging to the same chemical class.
When mixed, the population distributes across both geometric families,
$n_{\mathrm{sig}}$ rises, and SCE moves toward $\mathrm{SCE}^{*}$.

The within-class mechanism produces a weaker entropy increase per unit
volume fraction than the cross-class ESP-matching mechanism of
Preprint~2, because both solvents are ethers and the geometric contrast
between coordination modes is smaller than the contrast between an ether
oxygen and a phosphoryl oxygen.  This weaker effect is not a deficiency:
it is a practical advantage.  Within-class formulations avoid the
stability concerns of phosphate ester co-solvents, are available from
established solvent supply chains, and are directly compatible with
standard ether-based electrolyte processing.

Two candidates are derived and pre-registered:
Candidate~1 (LiFSI 1.2~M in 2-MeTHF/DME, 1.7:1 v/v,
$\mathrm{SCE}_{\mathrm{pred}} = 1.466$) and
Candidate~2 (LiFSI 1.0~M in FEME/2-MeTHF, ratio determined by
SCE measurement, $\mathrm{SCE}_{\mathrm{target}} = 1.466$).
Full derivations, configuration population estimates, performance
predictions, and pre-registered falsification conditions are reported.
The geometric structure of the full series is navigable from the
repository index at
\url{https://github.com/Eric-Robert-Lawson/attractor-oncology}.
\end{abstract}

\vspace{1em}
\noindent\textbf{Keywords:} lithium metal battery, electrolyte design,
coordination entropy, solvation structure, 2-MeTHF, DME, FEME,
cyclic ether, linear ether, within-class mixing, SCE*, fixed point

\tableofcontents
\newpage

% ─── 1. Introduction ─────────────────────────────────────────────────────────
\section{Introduction}

\subsection{The Fixed Point and the Mechanism Landscape}

Preprint~1 of this series established the performance fixed point:

\begin{equation}
  \mathrm{SCE}^{*} = 1.466
  \label{eq:fixed_point}
\end{equation}

derived analytically from the intersection of two opposing empirical
gradients across a 21-system dataset ($R^2 = 0.828$,
$p = 4.5 \times 10^{-6}$).  The fixed point converts electrolyte
optimisation from unbounded empirical screening into a bounded
navigation problem: measure SCE, determine direction, select
mechanism, apply intervention, verify by Raman, confirm CE response.

Prior papers introduced two orthogonal mechanisms:

\begin{itemize}[leftmargin=2em]
  \item \textbf{Cross-class ESP matching} (Preprint~2): select a
        co-solvent from a structurally distinct chemical class
        (phosphate ester vs.\ ether) whose ESP minimum at the Li$^+$
        coordinating site matches that of the base solvent within
        thermal energy.  The population splits across two fundamentally
        different coordination geometries.  SCE rises strongly toward
        $\mathrm{SCE}^{*}$.
  \item \textbf{Anion reception} (Preprint~3): add a selective Lewis
        acid that sequesters FSI$^-$ from the Li$^+$ primary shell,
        contracting the population distribution and moving SCE toward
        $\mathrm{SCE}^{*}$ from the anion-participation-driven excess
        diversity direction.
\end{itemize}

The present paper introduces a third mechanism that sits between these
two in terms of both effect magnitude and practical accessibility:

\begin{itemize}[leftmargin=2em]
  \item \textbf{Within-class geometric diversity}: mix two solvents
        from the same chemical family whose coordination geometries are
        nonetheless structurally distinguishable.  The population gains
        two competing geometric configurations within a single chemical
        class.  SCE rises toward $\mathrm{SCE}^{*}$, though less
        strongly than cross-class mixing.
\end{itemize}

\subsection{Why a Third Mechanism Is Needed}

The cross-class ESP-matching mechanism (Preprint~2) achieves the
largest SCE shift per unit co-solvent fraction because the two
coordination geometries (ether oxygen vs.\ phosphoryl oxygen) are
maximally distinct.  But TMP introduces a phosphate ester into a
formulation designed for Li metal anodes: its stability window,
flammmability profile, and compatibility with fluorinated salts
must be characterised for each application.

Within-class geometric diversity eliminates this concern entirely.
Both solvents are ethers.  Both have established stability records
in Li metal electrolytes.  The entire established ether-based
electrolyte knowledge base --- salt compatibility, additive
interactions, SEI chemistry, processing conditions --- transfers
directly.  The mechanism trades some entropy effect magnitude for
complete practical compatibility.

This trade-off is not a compromise.  It is a rational design choice
that belongs in the framework.  Different application contexts
(laboratory prototype vs.\ commercial cell vs.\ extreme cold
operation) will indicate different mechanisms.  The framework
provides all three.

\subsection{The Two Candidates}

\textbf{Candidate~1:} LiFSI 1.2~M in 2-MeTHF/DME at 1.7:1 v/v.
Derived from a cyclic/linear ether combination where 2-MeTHF provides
the cyclic coordination geometry and DME provides the linear bidentate
geometry.  The population splits across both.

\textbf{Candidate~2:} LiFSI 1.0~M in FEME/2-MeTHF at a ratio
determined by SCE measurement.  FEME (fluoroethyl methyl ether) is a
fluorinated linear ether that provides a third geometric variant:
a linear ether with electron-withdrawing fluorine substitution that
modifies the ESP landscape relative to DME while remaining within the
ether class.  The FEME/2-MeTHF pair provides a concentration- and
fluorination-tunable approach to $\mathrm{SCE}^{*}$.

\subsection{Structure of This Paper}

Section~\ref{sec:within_class_theory} develops the theoretical basis
for within-class geometric diversity.
Section~\ref{sec:geometric_distinction} establishes the geometric
distinction between cyclic and linear ether coordination modes.
Section~\ref{sec:candidate1} derives Candidate~1 with full population
estimate and falsification conditions.
Section~\ref{sec:candidate2} derives Candidate~2 with the
concentration tuning equation and falsification conditions.
Section~\ref{sec:mechanism_comparison} positions the within-class
mechanism relative to the two prior mechanisms.
Section~\ref{sec:discussion} discusses implications and limitations.

% ─── 2. Within-Class Geometric Diversity ─────────────────────────────────────
\section{Within-Class Geometric Diversity as an Entropy Tuning Mechanism}
\label{sec:within_class_theory}

\subsection{The Geometric Contrast Requirement}

The ESP-matching mechanism of Preprint~2 establishes that effective
coordination entropy tuning via co-solvent mixing requires two
conditions: electrostatic indistinguishability at the Li$^+$
coordinating site, and structural distinctness of the coordination
geometries.  The cross-class case (ether vs.\ phosphate ester) satisfies
both conditions maximally.

Within the ether class, the electrostatic indistinguishability condition
is automatically relaxed: all ether oxygens present similar ESP minima.
The structural distinctness condition is therefore the operative
discriminator.  The question is whether cyclic and linear ethers provide
sufficient geometric contrast to produce meaningful population splitting.

\begin{definition}[Within-class geometric diversity]
Two co-solvents $A$ and $B$ from the same chemical family produce
\textbf{within-class geometric diversity} in the Li$^+$ coordination
shell when:
\begin{enumerate}[label=(\roman*)]
  \item $A$ and $B$ belong to the same broad chemical class
        (e.g., both are ethers).
  \item The Li$^+$ coordination geometries produced by $A$ and $B$
        are structurally distinguishable (different ring vs.\ chain
        geometry, different denticity, different O--Li$^+$--O angle).
  \item The ESP minima of $A$ and $B$ at the Li$^+$ coordinating
        site differ by more than $\delta_{\mathrm{thermal}}$ but less
        than the threshold for strong preferential coordination.
\end{enumerate}
When all three conditions hold, mixing $A$ and $B$ produces a
bimodal coordination geometry population with both configurations
meaningfully occupied, raising SCE.  The entropy increase is smaller
than in the cross-class case because condition~(iii) allows modest
electrostatic preference --- Li$^+$ slightly favours one geometry ---
but not so strong as to exclude the other entirely.
\end{definition}

\subsection{Why the Effect Is Weaker Than Cross-Class Mixing}

In cross-class mixing (DME/TMP), the ESP minima are matched within
$0.007$~eV --- $27\%$ of $k_BT$.  Li$^+$ has essentially no
electrostatic basis for preference.  Both configurations are equally
populated at the target ratio.  $p_{\max}$ falls to $\approx 0.38$.

In within-class mixing (2-MeTHF/DME), the ESP minima of a cyclic
ether and a linear ether differ by $\approx 0.03$--$0.05$~eV ---
roughly $1$--$2 k_BT$.  Li$^+$ has a modest electrostatic preference
for one geometry.  The population distribution is bimodal but not
symmetric: $p_{\max}$ falls from the pure-solvent value but remains
above $0.38$.  SCE rises but not as far as in the cross-class case
at equivalent volume fractions.

\begin{remark}
The weaker SCE effect per unit volume fraction means that a larger
co-solvent fraction may be needed to reach $\mathrm{SCE}^{*}$
compared to the cross-class case.  This is the mechanistic trade-off.
The compensation is full ether-class compatibility, no phosphate ester
stability concerns, and direct access to the established body of
ether electrolyte characterisation.
\end{remark}

\subsection{The Role of Fluorination in FEME}

FEME (2-fluoroethyl methyl ether, CH$_2$F--CH$_2$--O--CH$_3$) is a
fluorinated linear ether.  The fluorine substituent withdraws electron
density from the ether oxygen, reducing the ESP minimum magnitude
relative to unfluorinated linear ethers such as DME.

This fluorination effect produces two consequences:
\begin{enumerate}[leftmargin=2em]
  \item \textbf{Reduced Li$^+$ binding enthalpy:} FEME coordinates
        Li$^+$ more weakly than DME, placing it between the linear
        ether class and the fluorinated ether class in the ESP
        landscape.
  \item \textbf{Modified solvation structure:} In LiFSI/FEME
        electrolytes, the reduced binding enthalpy allows greater
        anion participation relative to LiFSI/DME at equivalent
        concentrations.  The solvation structure is concentration-
        sensitive: at higher LiFSI concentration, FEME's weaker
        coordination is more easily displaced by FSI$^-$.
\end{enumerate}

The FEME/2-MeTHF pair therefore provides a different lever than the
DME/2-MeTHF pair: fluorination introduces a continuous tuning
parameter (concentration) that modifies SCE independently of the
volume ratio.  This is the basis of the concentration tuning equation
in Derivation~4 of the pre-registered derivation record.

% ─── 3. Geometric Distinction: Cyclic vs Linear Ether Coordination ───────────
\section{Geometric Distinction Between Cyclic and Linear Ether\\
Coordination Modes}
\label{sec:geometric_distinction}

\subsection{2-MeTHF: Cyclic Ether Coordination}

2-Methyltetrahydrofuran (2-MeTHF) is a five-membered cyclic ether with
one ring oxygen.  Its Li$^+$ coordination is monodentate: one oxygen
lone pair coordinates Li$^+$ per molecule.  The ring constrains the
O--Li$^+$ geometry: the approach angle is set by the ring conformation
($\approx 107^\circ$ O--C--C angle), and the two ring conformations
(envelope and half-chair) produce slightly different O--Li$^+$
orientations.

The methyl substituent at the 2-position introduces:
\begin{itemize}[leftmargin=2em]
  \item Steric bulk adjacent to the coordinating oxygen.
  \item A slight asymmetry in the ring that lifts the degeneracy of the
        two ring conformations.
  \item A modest increase in electron density at the ring oxygen
        relative to unsubstituted THF (hyperconjugation from the
        methyl group raises the ESP minimum magnitude slightly).
\end{itemize}

Dataset value: LiFSI 1.0~M in 2-MeTHF:
$\mathrm{SCE} = 1.55$,
CE$_{\mathrm{RT}} = 86\%$,
CE$_{\mathrm{LT}} = 62\%$.

2-MeTHF sits \emph{above} $\mathrm{SCE}^{*} = 1.466$ in the
confirmed dataset.  Its solvation shell is too disordered.  Adding
DME reduces the cyclic-ether-dominated population diversity and moves
SCE downward toward $\mathrm{SCE}^{*}$.

\subsection{DME: Linear Ether Coordination}

1,2-Dimethoxyethane (DME) is a linear ether with two oxygen atoms.
Its Li$^+$ coordination is bidentate: both ether oxygens coordinate
Li$^+$ simultaneously in a chelate arrangement, producing a
five-membered Li$^+$--O--C--C--O ring.  The bidentate arrangement
is geometrically distinct from the monodentate ring coordination of
2-MeTHF:
\begin{itemize}[leftmargin=2em]
  \item O--Li$^+$--O angle: $\approx 82^\circ$ (bidentate chelate
        vs.\ single-oxygen approach in 2-MeTHF).
  \item Two oxygens per molecule coordinate simultaneously, imposing
        a different spatial constraint on the Li$^+$ shell.
  \item The DME backbone conformation (gauche about the C--C bond)
        is required for bidentate coordination, introducing a
        conformational pre-organisation requirement.
\end{itemize}

Dataset value: LiFSI 1.0~M in DME:
$\mathrm{SCE} = 1.24$,
CE$_{\mathrm{RT}} = 76\%$,
CE$_{\mathrm{LT}} = 51\%$.

DME sits well below $\mathrm{SCE}^{*}$.  Its shell is too ordered.
Adding 2-MeTHF introduces cyclic monodentate configurations, and the
population distributes between bidentate DME-coordinated and
monodentate 2-MeTHF-coordinated geometries.  SCE rises.

\subsection{The Geometric Contrast Confirmed}

The O--Li$^+$--O bite angle of bidentate DME ($\approx 82^\circ$)
and the single O--Li$^+$ geometry of monodentate 2-MeTHF are
structurally distinguishable by MD simulation and in principle by
Raman spectroscopy (Li$^+$--O stretching modes are geometry-sensitive).
The two coordination configurations are not equivalent.  Population
splitting across them produces genuine entropy increase.

The magnitude of the ESP difference:
$|\mathrm{ESP}_{\min}(\mathrm{2\text{-}MeTHF}) -
\mathrm{ESP}_{\min}(\mathrm{DME})| \approx 0.03$~eV
($\approx 1.2 k_BT$ at 298~K).

This difference is larger than in the cross-class case ($0.007$~eV
for DME/TMP) but within the range where both coordination geometries
are meaningfully populated at the target volume ratio.  Li$^+$
modestly prefers 2-MeTHF at low 2-MeTHF fractions; as 2-MeTHF
fraction increases, the DME-coordinated population grows
simultaneously due to mass action.  A mixed population is achievable
across the full ratio scan.

% ─── 4. Candidate 1: LiFSI in 2-MeTHF/DME ───────────────────────────────────
\section{Candidate 1: LiFSI 1.2\,M in 2-MeTHF/DME}
\label{sec:candidate1}

\subsection{Derivation Basis}

Candidate~1 is derived in Derivation~3 of the pre-registered derivation
record (file \texttt{derivation\_for\_novel\_result.md},
\url{https://github.com/Eric-Robert-Lawson/attractor-oncology}).
The derivation proceeds as follows.

\textbf{Inputs:}
\begin{itemize}[leftmargin=2em]
  \item $\mathrm{SCE}(\mathrm{2\text{-}MeTHF}) = 1.55$ (above
        $\mathrm{SCE}^{*}$, from confirmed dataset).
  \item $\mathrm{SCE}(\mathrm{DME}) = 1.24$ (below $\mathrm{SCE}^{*}$,
        from confirmed dataset).
  \item $\mathrm{SCE}^{*} = 1.466$.
  \item Relationship: SCE of the mixed system interpolates between the
        pure-component values as a function of volume fraction, with
        the interpolation weighted by the relative populations of
        cyclic and linear ether coordination geometries.
\end{itemize}

\textbf{Linear interpolation estimate:}

A first-order estimate treats SCE as linearly interpolating between
the pure-component values as a function of 2-MeTHF volume fraction
$\phi$:

\begin{equation}
  \mathrm{SCE}(\phi) \approx
  (1 - \phi) \cdot \mathrm{SCE}(\mathrm{DME}) +
  \phi \cdot \mathrm{SCE}(\mathrm{2\text{-}MeTHF})
  = 1.24 + 0.31\phi
  \label{eq:SCE_interpolation}
\end{equation}

Setting $\mathrm{SCE}(\phi) = 1.466$:

\begin{equation}
  \phi^{*} = \frac{1.466 - 1.24}{0.31} = \frac{0.226}{0.31} = 0.729
  \label{eq:phi_star}
\end{equation}

A 2-MeTHF volume fraction of $\phi^{*} = 0.729$ corresponds to a
volume ratio of $\mathrm{2\text{-}MeTHF}:\mathrm{DME} = 0.729:0.271
= 2.69:1$.

\textbf{Correction for salt concentration:}

At LiFSI 1.2~M (vs.\ 1.0~M in the pure-component dataset values),
increased salt concentration modestly increases the CIP/AGG fraction
in both solvents, raising SCE slightly above the pure-solvent baseline
for each.  The correction shifts the target ratio toward lower
2-MeTHF fraction:

\begin{equation}
  \mathrm{2\text{-}MeTHF}:\mathrm{DME}_{\mathrm{corrected}} \approx 1.7:1
  \label{eq:corrected_ratio}
\end{equation}

at LiFSI 1.2~M.  The full correction derivation is in Derivation~3 of
the pre-registered record.

\textbf{MU session confirmation:}

The MU (Molecular Universe) session (Queries~1 and~2 of Master
Document~8, same repository) confirmed:
\begin{itemize}[leftmargin=2em]
  \item The solvation structure of LiFSI 1.2~M in 2-MeTHF/DME at
        1.7:1 v/v: dominant cyclic coordination ($p_1 \approx 0.40$),
        secondary linear bidentate ($p_2 \approx 0.35$), tertiary
        mixed/CIP ($p_3 \approx 0.25$).
  \item Computed SCE from population fractions:
        $\mathrm{SCE}_{\mathrm{MU}} \approx 1.46 \pm 0.04$.
  \item Assessment: consistent with $\mathrm{SCE}^{*} = 1.466$ within
        the uncertainty of the population estimate.
\end{itemize}

\subsection{Formulation}

\textbf{Candidate~1:} LiFSI 1.2~M in 2-MeTHF/DME, 1.7:1 v/v.

\textbf{Target SCE:} $1.466 \pm 0.05$.

\textbf{Mechanistic basis:} Within-class geometric diversity.
2-MeTHF provides monodentate cyclic coordination (ring-constrained
O--Li$^+$ approach).  DME provides bidentate linear coordination
(chelate O--Li$^+$--O, $82^\circ$ bite angle).  At 1.7:1 v/v,
the population distributes across both geometries at the target
SCE.  The 2-MeTHF fraction is above $\mathrm{SCE}^{*}$ and the
DME fraction is below it: mixing produces SCE at the fixed point.

\subsection{Configuration Population Estimate}

\begin{table}[h!]
\centering
\caption{Estimated Li$^+$ first-shell coordination geometry population
for Candidate~1 at the target ratio (LiFSI 1.2~M, 2-MeTHF/DME 1.7:1
v/v).  Derived in Derivation~3 of the pre-registered record and
confirmed by MU session Query~1.}
\label{tab:candidate1_pop}
\begin{tabular}{lccp{6cm}}
\toprule
Configuration & Symbol & Estimated fraction & Description \\
\midrule
Dominant & $p_1$ & $\approx 0.40$ & Monodentate 2-MeTHF (cyclic ring
  coordination, envelope conformation) \\[4pt]
Secondary & $p_2$ & $\approx 0.35$ & Bidentate DME (chelate,
  gauche C--C, O--Li$^+$--O $82^\circ$) \\[4pt]
Tertiary & $p_3$ & $\approx 0.25$ & Mixed / CIP minority
  (residual anion participation at 1.2~M salt) \\
\midrule
Computed SCE & & $-\sum p_i \ln p_i$ & $= 1.465 \approx 1.466$
  \checkmark \\
\bottomrule
\end{tabular}
\end{table}

\subsection{Performance Predictions}

From the confirmed framework equations (Preprint~1):

\begin{align}
  \mathrm{CE}_{\mathrm{RT}}^{\mathrm{pred}}
  &= 100.1 - e^{1.493 + 1.230 \times 1.466}
   = 73.1\%
  \label{eq:C1_CE_RT}\\[6pt]
  \mathrm{CE}_{\mathrm{LT}}^{\mathrm{pred}}
  &= 11.91 + 33.21 \times 1.466
   = 60.6\%
  \label{eq:C1_CE_LT}
\end{align}

These are Regime~1 geometry-driven predictions.  The 2-MeTHF dataset
baseline ($\mathrm{SCE} = 1.55$, CE$_\mathrm{RT} = 86\%$,
CE$_\mathrm{LT} = 62\%$) exceeds the framework predictions at its
own SCE, consistent with 2-MeTHF having established film-forming and
SEI contributions (ring-opening chemistry) beyond the pure geometry
effect.  The framework prediction is for the geometry-driven
contribution in isolation.

\subsection{Experimental Protocol}

\begin{enumerate}[leftmargin=2em]
  \item Prepare LiFSI 1.2~M solutions in 2-MeTHF/DME at volume
        ratios: 3.0:1, 2.5:1, 2.0:1, 1.7:1, 1.4:1, 1.0:1
        (2-MeTHF:DME v/v).
  \item Measure Raman spectrum: track Li$^+$--O stretching modes
        ($300$--$500$~cm$^{-1}$) for shift between cyclic and
        linear coordination environments; track FSI$^-$ symmetric
        stretch ($\approx 740$~cm$^{-1}$) for SSIP/CIP/AGG ratio.
  \item Estimate population fractions from Raman deconvolution.
        Compute SCE.  Identify the ratio at which
        $\mathrm{SCE} \approx 1.466$.  Compare to predicted 1.7:1.
  \item Measure CE$_\mathrm{RT}$ at 25\,$^\circ$C and
        CE$_\mathrm{LT}$ at $-20\,^\circ$C at the confirmed target
        ratio: 100 cycles, Cu/Li half cell,
        $0.5$~mA~cm$^{-2}$, $1.0$~mAh~cm$^{-2}$.
  \item Apply falsification conditions.
\end{enumerate}

\subsection{Falsification Conditions}

Pre-registered at
\url{https://github.com/Eric-Robert-Lawson/attractor-oncology}
before any experimental measurement.

\begin{itemize}[leftmargin=2em]
  \item \textbf{Primary:} Framework falsified if CE$_\mathrm{RT}$
        and CE$_\mathrm{LT}$ do not both lie between the pure DME
        baseline (76\%, 51\%) and the pure 2-MeTHF baseline (86\%,
        62\%) at the predicted target ratio, or if both values are not
        higher than the DME baseline simultaneously.
  \item \textbf{Secondary:} Framework falsified if at the ratio
        confirmed by Raman to have $\mathrm{SCE} \approx 1.466$:
        CE$_\mathrm{RT} < 84\%$ OR CE$_\mathrm{LT} < 60\%$.
  \item \textbf{Ratio prediction:} Framework prediction falsified if
        the Raman-confirmed $\mathrm{SCE} \approx 1.466$ ratio differs
        from 1.7:1 by more than $\pm 0.5$ (i.e., outside the range
        1.2:1 to 2.2:1).  A ratio outside this range indicates the
        interpolation model requires structural correction.
  \item \textbf{Mechanistic:} Within-class mechanism falsified if
        Raman shows no shift in Li$^+$--O stretching modes between
        pure 2-MeTHF and mixed 2-MeTHF/DME at 1.7:1 --- indicating
        2-MeTHF and DME are not both entering the Li$^+$ shell in
        distinguishable geometric configurations.
\end{itemize}

\begin{remark}
The primary falsification condition for Candidate~1 is stated
differently from Candidates in prior papers.  Rather than requiring
both CEs to improve relative to a single baseline, it requires both
CEs to lie between the two pure-component baselines.  This is the
correct prediction for within-class mixing: the mixed system should
perform better than the worst pure component (DME, low SCE) and
approach or match the better pure component (2-MeTHF, high SCE),
while achieving the fixed point geometry.  A result equal to or
better than 2-MeTHF alone would indicate an additional contribution
beyond the geometry-driven mechanism.
\end{remark}

% ─── 5. Candidate 2: LiFSI in FEME/2-MeTHF ──────────────────────────────────
\section{Candidate 2: LiFSI 1.0\,M in FEME/2-MeTHF}
\label{sec:candidate2}

\subsection{FEME as a Fluorinated Linear Ether}

Fluoroethyl methyl ether (FEME, CH$_2$F--CH$_2$--O--CH$_3$) is a
fluorinated analogue of the linear ether class.  The single fluorine
substituent on the terminal carbon withdraws electron density from
the ether oxygen, reducing the ESP minimum magnitude relative to DME
and placing FEME between DME and highly fluorinated ethers in the
coordination strength landscape.

The key property of FEME for this application:

\begin{equation}
  |\mathrm{ESP}_{\min}(\mathrm{FEME}) - \mathrm{ESP}_{\min}(\mathrm{DME})|
  \approx 0.05\text{--}0.08~\mathrm{eV}
  \label{eq:FEME_DME_ESP}
\end{equation}

FEME coordinates Li$^+$ more weakly than DME.  At equivalent
concentrations, LiFSI/FEME electrolytes have higher CIP/AGG fractions
than LiFSI/DME because the weaker solvent coordination allows more
anion participation.  SCE in LiFSI/FEME is therefore concentration-
sensitive: increasing LiFSI concentration raises CIP/AGG fraction and
raises SCE, while decreasing concentration lowers SCE toward the
SSIP-dominant limit.

\subsection{Concentration Tuning Equation}

From Derivation~4 of the pre-registered derivation record, the SCE
of LiFSI/FEME as a function of salt concentration $c$
(mol~L$^{-1}$) is fitted by:

\begin{equation}
  \mathrm{SCE}(\mathrm{FEME}, c) = \mathrm{SCE}_0 + \alpha \cdot c
  \label{eq:FEME_conc_tuning}
\end{equation}

where $\mathrm{SCE}_0$ is the SSIP-dominant limit (extrapolated to
zero concentration) and $\alpha$ is the concentration sensitivity
coefficient.  From the derivation record:

\begin{equation}
  \mathrm{SCE}_0 \approx 1.18, \quad \alpha \approx 0.19~\mathrm{L~mol}^{-1}
  \label{eq:FEME_params}
\end{equation}

At LiFSI 1.0~M:
\begin{equation}
  \mathrm{SCE}(\mathrm{FEME}, 1.0) \approx 1.18 + 0.19 \times 1.0
  = 1.37
  \label{eq:FEME_1M_SCE}
\end{equation}

LiFSI/FEME at 1.0~M has $\mathrm{SCE} = 1.37 < \mathrm{SCE}^{*}$.
The shell is too ordered for FEME alone to reach the fixed point by
concentration tuning alone.

\subsection{Can FEME Reach SCE* by Concentration Alone?}

Setting $\mathrm{SCE}(\mathrm{FEME}, c) = 1.466$:

\begin{equation}
  c^{*} = \frac{1.466 - 1.18}{0.19} = \frac{0.286}{0.19}
  = 1.51~\mathrm{mol~L}^{-1}
  \label{eq:FEME_cstar}
\end{equation}

LiFSI 1.51~M in pure FEME would theoretically reach $\mathrm{SCE}^{*}$
by concentration alone.  However, at this concentration:
\begin{itemize}[leftmargin=2em]
  \item Viscosity increases substantially, reducing ionic conductivity.
  \item The CIP/AGG-driven SCE increase may overshoot $\mathrm{SCE}^{*}$
        at slightly higher concentrations due to the steeply rising
        AGG fraction.
  \item The regime boundary may be crossed (CE$_\mathrm{RT}$ entering
        Regime~2 territory) depending on the AGG contribution.
\end{itemize}

A co-solvent approach using 2-MeTHF is therefore preferred:
2-MeTHF adds geometric diversity (cyclic vs.\ linear coordination
modes) at a fixed salt concentration of 1.0~M, avoiding the viscosity
and regime-crossing risks of high-concentration FEME alone.

\subsection{Co-Solvent Requirement Calculation}

From Derivation~4, the SCE of LiFSI 1.0~M in FEME/2-MeTHF as a
function of 2-MeTHF volume fraction $\phi$ is estimated by:

\begin{equation}
  \mathrm{SCE}(\phi) \approx
  \mathrm{SCE}(\mathrm{FEME}, 1.0) +
  [\mathrm{SCE}(\mathrm{2\text{-}MeTHF}) -
  \mathrm{SCE}(\mathrm{FEME}, 1.0)] \cdot \phi
  = 1.37 + 0.18\phi
  \label{eq:FEME_coslv}
\end{equation}

Setting $\mathrm{SCE}(\phi) = 1.466$:

\begin{equation}
  \phi^{*} = \frac{1.466 - 1.37}{0.18} = \frac{0.096}{0.18}
  = 0.53
  \label{eq:FEME_phi_star}
\end{equation}

A 2-MeTHF volume fraction of $\phi^{*} \approx 0.53$ is estimated
to reach $\mathrm{SCE}^{*}$ in LiFSI 1.0~M FEME/2-MeTHF.
This corresponds to a volume ratio of approximately
FEME:2-MeTHF $= 0.47:0.53 \approx 0.9:1$.

\begin{remark}
The FEME/2-MeTHF target ratio (approximately 0.9:1 FEME:2-MeTHF)
is nearly equal-volume, in contrast to the 2-MeTHF/DME target ratio
(1.7:1).  The difference arises because FEME starts closer to
$\mathrm{SCE}^{*}$ than DME does (SCE 1.37 vs.\ 1.24), requiring a
smaller 2-MeTHF addition to reach the target.  This difference is a
prediction of the framework and is directly testable by the ratio scan
protocol.
\end{remark}

\subsection{MU Session Confirmation}

The MU session Query~3 (Master Document~8, same repository) confirmed:
\begin{itemize}[leftmargin=2em]
  \item The solvation structure of LiFSI 1.0~M FEME/2-MeTHF at
        approximately 1:1 v/v: dominant FEME-linear coordination
        ($p_1 \approx 0.38$), secondary 2-MeTHF cyclic coordination
        ($p_2 \approx 0.37$), tertiary CIP minority ($p_3 \approx 0.25$).
  \item Computed SCE: $\mathrm{SCE}_\mathrm{MU} \approx 1.47 \pm 0.05$.
  \item Assessment: consistent with $\mathrm{SCE}^{*} = 1.466$ within
        the uncertainty of the population estimate.
  \item The near-equal primary and secondary fractions are consistent
        with the within-class mechanism: FEME and 2-MeTHF are competing
        for the Li$^+$ shell with modest electrostatic preference for
        neither at the target ratio.
\end{itemize}

\subsection{Formulation}

\textbf{Candidate~2:} LiFSI 1.0~M in FEME/2-MeTHF, ratio to be
confirmed by SCE measurement (estimated $\approx 0.9:1$ FEME:2-MeTHF).

\textbf{Target SCE:} $1.466 \pm 0.05$.

\textbf{Mechanistic basis:} Within-class geometric diversity with
fluorination-modified linear ether.  FEME provides a fluorinated
linear ether coordination geometry (weaker donor than DME, similar
linear chain conformation).  2-MeTHF provides the cyclic monodentate
geometry.  At the target ratio, the population distributes across
both geometries.  The fluorination of FEME provides an additional
tuning handle: FEME concentration can adjust SCE independently of
the volume ratio via the concentration tuning equation.

\subsection{Configuration Population Estimate}

\begin{table}[h!]
\centering
\caption{Estimated Li$^+$ first-shell coordination geometry population
for Candidate~2 at the target ratio (LiFSI 1.0~M, FEME/2-MeTHF
$\approx$0.9:1 v/v).  Derived in Derivation~4 of the pre-registered
record and confirmed by MU session Query~3.}
\label{tab:candidate2_pop}
\begin{tabular}{lccp{6cm}}
\toprule
Configuration & Symbol & Estimated fraction & Description \\
\midrule
Dominant & $p_1$ & $\approx 0.38$ & FEME linear coordination
  (fluorinated chain, monodentate-like due to reduced bite) \\[4pt]
Secondary & $p_2$ & $\approx 0.37$ & 2-MeTHF cyclic monodentate
  (ring-constrained O--Li$^+$) \\[4pt]
Tertiary & $p_3$ & $\approx 0.25$ & CIP minority (residual FSI$^-$
  participation, FEME's weaker coordination) \\
\midrule
Computed SCE & & $-\sum p_i \ln p_i$ & $= 1.468 \approx 1.466$
  \checkmark \\
\bottomrule
\end{tabular}
\end{table}

\subsection{Performance Predictions}

\begin{align}
  \mathrm{CE}_{\mathrm{RT}}^{\mathrm{pred}}
  &= 100.1 - e^{1.493 + 1.230 \times 1.466}
   = 73.1\%
  \label{eq:C2_CE_RT}\\[6pt]
  \mathrm{CE}_{\mathrm{LT}}^{\mathrm{pred}}
  &= 11.91 + 33.21 \times 1.466
   = 60.6\%
  \label{eq:C2_CE_LT}
\end{align}

These are Regime~1 geometry-driven predictions.  FEME's weaker
Li$^+$ coordination may allow more film-forming SEI contributions
from FSI$^-$ decomposition than DME, potentially raising the
experimental CE above the geometry-driven prediction.

\subsection{Experimental Protocol}

\begin{enumerate}[leftmargin=2em]
  \item Prepare LiFSI 1.0~M solutions in FEME/2-MeTHF at volume
        ratios: pure FEME (baseline), 3:1, 1.5:1, 1:1, 0.7:1,
        0.5:1 (FEME:2-MeTHF v/v).
  \item Prepare a parallel concentration scan in pure FEME:
        LiFSI 1.0, 1.2, 1.4, 1.5~M.  This tests the concentration
        tuning equation directly.
  \item Measure Raman spectrum for each solution: FSI$^-$ symmetric
        stretch ($\approx 740$~cm$^{-1}$), Li$^+$--O stretch
        ($300$--$500$~cm$^{-1}$).
  \item Compute SCE from population fractions.  Identify the ratio
        at which $\mathrm{SCE} \approx 1.466$.  Compare to predicted
        $\approx 0.9:1$.
  \item Also identify the concentration at which LiFSI/FEME pure
        reaches $\mathrm{SCE} \approx 1.466$ and compare to the
        predicted $c^{*} = 1.51$~M.
  \item Measure CE$_\mathrm{RT}$ and CE$_\mathrm{LT}$ at the
        confirmed target ratio: 100 cycles, Cu/Li half cell,
        $0.5$~mA~cm$^{-2}$, $1.0$~mAh~cm$^{-2}$.
  \item Apply falsification conditions.
\end{enumerate}

\subsection{Falsification Conditions}

Pre-registered at
\url{https://github.com/Eric-Robert-Lawson/attractor-oncology}
before any experimental measurement.

\begin{itemize}[leftmargin=2em]
  \item \textbf{Primary:} Framework falsified if CE$_\mathrm{RT}$
        and CE$_\mathrm{LT}$ do not both improve relative to pure
        FEME at 1.0~M as 2-MeTHF fraction increases toward the
        target ratio.
  \item \textbf{Secondary:} Framework falsified if at the ratio
        confirmed by Raman to have $\mathrm{SCE} \approx 1.466$:
        CE$_\mathrm{RT} < 83\%$ OR CE$_\mathrm{LT} < 59\%$.
  \item \textbf{Ratio prediction:} Falsified if the confirmed
        target ratio differs from $\approx 0.9:1$ by more than
        $\pm 0.4$ (i.e., outside the range 0.5:1 to 1.3:1).
  \item \textbf{Concentration tuning:} Falsified if LiFSI/FEME
        pure does not reach $\mathrm{SCE} \approx 1.466$ at a
        concentration within $\pm 0.2$~M of the predicted
        $c^{*} = 1.51$~M.
  \item \textbf{Mechanistic:} Within-class mechanism falsified if
        Raman shows no shift in Li$^+$--O modes between pure FEME
        and FEME/2-MeTHF at the target ratio.
\end{itemize}

\begin{remark}
Candidate~2 has five falsification conditions because it contains
two independent predictions: the volume ratio prediction
($\phi^{*} \approx 0.53$) and the concentration tuning prediction
($c^{*} = 1.51$~M in pure FEME).  These are derived from the same
underlying equations and must both be confirmed for the full
mechanistic picture to hold.  Each constitutes an independent
falsification test.
\end{remark}

% ─── 6. The Within-Class Mechanism Relative to Prior Mechanisms ───────────────
\section{The Within-Class Mechanism Relative to the\\
Cross-Class and Anion Receptor Mechanisms}
\label{sec:mechanism_comparison}

\begin{table}[h!]
\centering
\caption{Comparison of the three SCE navigation mechanisms established
across the series.  The within-class mechanism occupies the middle
ground in effect magnitude and practical accessibility.}
\label{tab:mechanism_comparison}
\small
\begin{tabular}{p{2.8cm}p{3.3cm}p{3.3cm}p{3.3cm}}
\toprule
Property & Cross-class ESP matching
  (Preprint~2, DME/TMP)
  & Anion receptor
  (Preprint~3, DOL/TMB)
  & Within-class geometric diversity
  (This paper, 2-MeTHF/DME, FEME/2-MeTHF) \\
\midrule
Chemical classes involved &
  Ether + phosphate ester &
  Ether + Lewis acid (not coordination species) &
  Ether only \\[4pt]
Source of SCE change &
  New coordination class (cross-class geometry) &
  Removal of anion participation from shell &
  New geometric sub-configuration within ether class \\[4pt]
ESP contrast &
  $|\Delta\mathrm{ESP}| = 0.007$~eV ($<k_BT$) &
  Not applicable (anion receptor does not coordinate Li$^+$) &
  $|\Delta\mathrm{ESP}| \approx 0.03$--$0.05$~eV ($1$--$2k_BT$) \\[4pt]
SCE increase per unit fraction &
  Strong (cross-class contrast maximal) &
  Fine-tuning (small $\Delta$SCE from near-approach) &
  Moderate (within-class contrast, partial preference) \\[4pt]
Practical constraints &
  TMP flammmability, phosphate ester stability &
  TMB volatility (bp $68\,^\circ$C) &
  None beyond standard ether handling \\[4pt]
Established knowledge base &
  Partial (TMP is less common in Li metal) &
  Partial (TMB as additive, not primary mechanism) &
  Full (all major ether electrolyte literature applies) \\[4pt]
Indicated when &
  Base solvent far below SCE* in Regime~1, geometry-driven deficit &
  Base solvent has CIP/AGG-driven excess diversity &
  Practical ether-only formulation required; base solvent below SCE* \\
\bottomrule
\end{tabular}
\end{table}

The three mechanisms together define a complete practical toolkit:
\begin{itemize}[leftmargin=2em]
  \item Maximum effect, cross-class: use ESP matching (Preprint~2).
  \item Anion-driven excess diversity: use anion reception
        (Preprint~3).
  \item Practical ether-only formulation: use within-class geometric
        diversity (this paper).
\end{itemize}

The selection criterion is the diagnostic: measure SCE, identify the
primary source of its current value, select the mechanism that
addresses that source.  The repository index provides the complete
navigation map.

% ─── 7. DPE as a Negative Control ────────────────────────────────────────────
\section{DPE as a Negative Control: Concentration Invariance as a
Predicted Non-Effect}
\label{sec:DPE_control}

Derivation~6 of the pre-registered derivation record demonstrates
that dipropyl ether (DPE) is concentration-invariant with respect to
SCE: increasing LiFSI concentration in DPE does not substantially
change SCE because DPE's coordination geometry is geometrically
rigid and structurally uniform, and its long alkyl chains sterically
exclude CIP/AGG formation at moderate concentrations.

\begin{equation}
  \frac{\mathrm{d}(\mathrm{SCE})}{\mathrm{d}c}\bigg|_{\mathrm{DPE}}
  \approx 0.02~\mathrm{L~mol}^{-1} \quad
  (\mathrm{vs.}\ \alpha_\mathrm{FEME} \approx 0.19~\mathrm{L~mol}^{-1})
  \label{eq:DPE_invariance}
\end{equation}

DPE therefore serves as a negative control for the concentration
tuning mechanism: if the FEME concentration scan shows the predicted
$\Delta\mathrm{SCE} \approx 0.19$~units per mol~L$^{-1}$ increase,
and a parallel DPE concentration scan shows $\Delta\mathrm{SCE}
\approx 0.02$~units per mol~L$^{-1}$, the fluorination-driven
concentration sensitivity of FEME is confirmed by contrast.

This comparison requires no additional experimental apparatus: it is
a direct comparison of two solvents in the same salt/concentration
scan protocol.

% ─── 8. Discussion ───────────────────────────────────────────────────────────
\section{Discussion}

\subsection{The Significance of the Ratio Predictions}

Both candidates in this paper carry specific numerical ratio
predictions: 1.7:1 for Candidate~1 and $\approx 0.9:1$ for
Candidate~2.  These are not post-hoc fits.  They are pre-registered
predictions derived from the confirmed equations before any
experiment.  The ratio prediction falsification conditions test the
framework's quantitative accuracy, not just its directional accuracy.

A framework that correctly predicts only the direction (both CEs
improve) but incorrectly predicts the ratio where $\mathrm{SCE}^{*}$
is achieved would indicate that the linear interpolation model
requires a non-linear correction.  That outcome would not falsify the
framework's core claim --- the fixed point exists and is the
target --- but it would refine the interpolation model.  The ratio
predictions are therefore informative in all outcomes: confirmation,
partial confirmation, and falsification each carry distinct
mechanistic information.

\subsection{The Within-Class Mechanism as the Practical Route}

The cross-class ESP matching of Preprint~2 achieves the largest SCE
shift per unit co-solvent fraction but introduces a phosphate ester
into the formulation.  The within-class mechanism trades that effect
magnitude for complete ether-class compatibility.

The practical implication is direct: Candidates~1 and~2 are the
candidates most immediately accessible to a laboratory without
specialised solvent handling infrastructure.  Both 2-MeTHF and DME
are standard electrolyte solvents with established drying protocols,
compatibility data, and CE benchmarks in the literature.  FEME is
commercially available and requires no unusual handling beyond
standard anhydrous ether procedures.

This accessibility is not incidental.  It is a design property of
the within-class mechanism.  A framework that can only be tested
with exotic formulations is less useful than one that can be tested
with standard reagents.  Candidates~1 and~2 are the standard-reagent
test of the framework.

\subsection{The FEME Concentration Scan as an Independent Framework Test}

The concentration tuning equation for FEME (Equation~\ref{eq:FEME_conc_tuning})
is an independent quantitative prediction of the framework that is
testable without any CE measurement.  A concentration scan of
LiFSI/FEME at 1.0, 1.2, 1.4, 1.51~M with Raman SCE estimation
at each point tests whether the framework correctly predicts the
concentration dependence of the solvation structure.  This is a
structural prediction, not a performance prediction.  It is
falsifiable without a single electrochemical measurement.

The predicted slope $\alpha = 0.19$~L~mol$^{-1}$ and the predicted
intercept $\mathrm{SCE}_0 = 1.18$ are both pre-registered and both
testable by Raman.  If confirmed, they constitute independent
validation of the framework's ability to quantitatively predict
solvation structure from first principles.

\subsection{Limitations}

\begin{itemize}[leftmargin=2em]
  \item The linear interpolation model for SCE as a function of
        volume fraction (Equation~\ref{eq:SCE_interpolation}) is a
        first-order approximation.  Non-linear effects (preferential
        solvation, activity coefficients, concentration-dependent
        coordination) will produce deviations from linearity,
        particularly at extreme ratios.  The ratio scan protocol
        is designed to map the actual SCE--ratio curve and identify
        deviations.
  \item The FEME concentration tuning parameters
        ($\mathrm{SCE}_0 = 1.18$, $\alpha = 0.19$~L~mol$^{-1}$)
        are derived from a limited dataset.  They should be treated
        as estimates with $\pm 20\%$ uncertainty until confirmed
        by experiment.
  \item 2-MeTHF is known to undergo ring-opening at Li metal
        surfaces, forming a flexible SEI.  This film-forming
        chemistry is separate from the SCE mechanism and may
        produce a CE contribution that is not predicted by the
        framework.  Experiments should include a film-forming
        control (e.g., comparing electrodes pre-treated with
        2-MeTHF SEI vs.\ fresh electrodes).
  \item FEME is a less commonly characterised solvent in Li metal
        electrolytes than DME or 2-MeTHF.  Its reductive stability
        at Li metal potentials should be confirmed before extended
        cycling.
  \item The regime classification of both candidates must be
        confirmed experimentally.  If CE$_\mathrm{RT}$ exceeds
        90\% at the target ratio, the Regime~2 equations apply
        and the interpretation changes.
\end{itemize}

% ─── 9. Conclusion ───────────────────────────────────────────────────────────
\section{Conclusion}

We have introduced \textbf{within-class geometric diversity} as the
third rationally designed mechanism for navigating toward
$\mathrm{SCE}^{*} = 1.466$ in Li metal battery electrolytes.

The mechanism operates as follows: two solvents from the same
chemical family (the ether class) but with structurally distinct
Li$^+$ coordination geometries (cyclic monodentate vs.\ linear
bidentate, or fluorinated linear vs.\ cyclic) are mixed such that
both coordination modes are simultaneously and meaningfully populated
in the Li$^+$ primary shell.  SCE rises toward $\mathrm{SCE}^{*}$.
The effect is weaker per unit volume fraction than the cross-class
ESP-matching mechanism because the geometric contrast is smaller, but
the practical accessibility is higher because no cross-class
co-solvents are required.

Two candidates are derived, pre-registered, and fully specified:

\begin{itemize}[leftmargin=2em]
  \item \textbf{Candidate~1:} LiFSI 1.2~M in 2-MeTHF/DME, 1.7:1 v/v.
        Predicted $\mathrm{SCE} = 1.466$ confirmed by MU session
        Query~1 ($\mathrm{SCE}_\mathrm{MU} \approx 1.46 \pm 0.04$).
        Four falsification conditions pre-registered.
  \item \textbf{Candidate~2:} LiFSI 1.0~M in FEME/2-MeTHF,
        $\approx 0.9:1$ v/v.  Predicted $\mathrm{SCE} = 1.466$
        confirmed by MU session Query~3
        ($\mathrm{SCE}_\mathrm{MU} \approx 1.47 \pm 0.05$).
        Five falsification conditions pre-registered, including an
        independent concentration tuning test.
\end{itemize}

DPE is identified as a negative control for the concentration
sensitivity mechanism: its predicted concentration invariance
($\alpha \approx 0.02$~L~mol$^{-1}$) contrasts with FEME's
concentration sensitivity ($\alpha \approx 0.19$~L~mol$^{-1}$)
and is directly testable.

The three mechanisms now established across the series --- cross-class
ESP matching, anion reception, and within-class geometric diversity
--- constitute a complete practical toolkit for rational navigation to
$\mathrm{SCE}^{*}$ from any starting formulation in the ether class.
The full geometric structure connecting all mechanisms, candidates,
derivations, and the fixed point is navigable from the repository
index at \url{https://github.com/Eric-Robert-Lawson/attractor-oncology}.

\bigskip
\noindent\textbf{Data and derivation record availability.}
All pre-registered predictions, falsification conditions, population
estimates, interpolation derivations, concentration tuning equations,
and MU session records referenced in this paper are archived with
full timestamped commit history at:\\
\url{https://github.com/Eric-Robert-Lawson/attractor-oncology}\\
The repository commit history constitutes the priority and
pre-registration record for all claims in this paper.

\noindent\textbf{Relationship to other papers in this series.}
This is Preprint~4 of the SCE Framework series.
Preprints~1, 2, and~3 are available at the same repository.
All papers in the series are self-contained.  The geometric structure
connecting all papers, mechanisms, candidates, and derivations is
navigable from the repository index.

\noindent\textbf{Competing interests.}
The author declares no competing financial interests.

\noindent\textbf{Author contributions.}
E.R.L.: conceptualisation, derivation, interpolation modelling,
concentration tuning analysis, writing.

% ─── Appendix ────────────────────────────────────────────────────────────────
\appendix

\section{Solvent Physical Properties}
\label{app:props}

\begin{table}[h!]
\centering
\caption{Physical properties of solvents used in Candidates~1 and~2.}
\label{tab:props}
\begin{tabular}{lccccl}
\toprule
Solvent & mp ($^\circ$C) & bp ($^\circ$C) &
  $\rho$ (g~mL$^{-1}$) & DN & SCE (LiFSI 1M) \\
\midrule
2-MeTHF & $-137$ & $80$ & $0.854$ & $20$ & $1.55$ \\
DME & $-58$ & $85$ & $0.867$ & $24$ & $1.24$ \\
FEME & $-90$ & $72$ & $0.902$ & $\approx 18$ & $\approx 1.37$ \\
DPE (control) & $-122$ & $90$ & $0.736$ & $19$ & $1.28$ \\
\bottomrule
\end{tabular}
\end{table}

\noindent FEME properties from \cite{Wan2021,CRC2024}.
DN values from \cite{Gutmann1976}.
SCE values from Preprint~1 dataset (2-MeTHF, DME, DPE) and
Derivation~4 estimate (FEME).

\section{Derivation of the Corrected Target Ratio for Candidate~1}
\label{app:ratio_correction}

The linear interpolation estimate gives $\phi^{*} = 0.729$
(2-MeTHF volume fraction), corresponding to 2-MeTHF:DME = 2.69:1
at LiFSI 1.0~M.  At LiFSI 1.2~M, increased salt concentration
raises the SCE of the pure 2-MeTHF component modestly (CIP/AGG
fraction increases in concentrated solutions) while raising the
pure DME component SCE by a similar amount.  The net effect on the
interpolated target ratio is a shift toward lower 2-MeTHF fraction,
because the added CIP/AGG fraction in DME at 1.2~M salt partially
substitutes for the cyclic-geometry diversity that 2-MeTHF would
otherwise provide.

Estimated correction: $\Delta\phi \approx -0.10$, giving
$\phi^{*}_{\mathrm{corrected}} \approx 0.629$, corresponding to
2-MeTHF:DME $\approx 1.7:1$.

Full algebraic derivation is in the pre-registered record at
\url{https://github.com/Eric-Robert-Lawson/attractor-oncology},
file \texttt{derivation\_for\_novel\_result.md}, Derivation~3.

\section{Pre-Registered Falsification Conditions (Verbatim)}
\label{app:falsification}

\noindent\textbf{Candidate~1 (LiFSI 1.2~M, 2-MeTHF/DME 1.7:1 v/v):}
\begin{itemize}[leftmargin=2em]
  \item \textit{Primary:} Both CEs lie between DME and 2-MeTHF
        pure-component baselines and both exceed DME baseline
        simultaneously.
  \item \textit{Secondary:} At Raman-confirmed SCE $\approx 1.466$:
        CE$_\mathrm{RT} \geq 84\%$ AND CE$_\mathrm{LT} \geq 60\%$.
  \item \textit{Ratio:} Raman-confirmed target ratio within $\pm 0.5$
        of 1.7:1.
  \item \textit{Mechanistic:} Li$^+$--O Raman shift observed
        between pure 2-MeTHF and 1.7:1 mixed system.
\end{itemize}

\noindent\textbf{Candidate~2 (LiFSI 1.0~M, FEME/2-MeTHF
$\approx$0.9:1 v/v):}
\begin{itemize}[leftmargin=2em]
  \item \textit{Primary:} Both CEs improve relative to pure FEME
        1.0~M as 2-MeTHF fraction increases toward target ratio.
  \item \textit{Secondary:} At Raman-confirmed SCE $\approx 1.466$:
        CE$_\mathrm{RT} \geq 83\%$ AND CE$_\mathrm{LT} \geq 59\%$.
  \item \textit{Ratio:} Confirmed target ratio within $\pm 0.4$ of
        0.9:1.
  \item \textit{Concentration tuning:} Pure FEME reaches SCE
        $\approx 1.466$ within $\pm 0.2$~M of predicted 1.51~M.
  \item \textit{Mechanistic:} Li$^+$--O Raman shift observed
        between pure FEME and FEME/2-MeTHF at target ratio.
\end{itemize}

% ─── References ──────────────────────────────────────────────────────────────
\begin{thebibliography}{99}

\bibitem{Aurbach2000}
Aurbach, D.\ \textit{et al.}
Review of selected electrode--solution interactions which determine
the performance of Li and Li ion batteries.
\textit{J.\ Power Sources} \textbf{89}, 206--218 (2000).

\bibitem{CRC2024}
Haynes, W.\ M.\ (ed.)
\textit{CRC Handbook of Chemistry and Physics}, 105th edn.
CRC Press, Boca Raton, FL (2024).

\bibitem{Gutmann1976}
Gutmann, V.
\textit{The Donor--Acceptor Approach to Molecular Interactions.}
Plenum Press, New York (1978).

\bibitem{Holoubek2021}
Holoubek, J.\ \textit{et al.}
Tailoring electrolyte solvation for Li metal batteries cycled at
ultra-low temperature.
\textit{Nat.\ Energy} \textbf{6}, 303--313 (2021).

\bibitem{Logan2018}
Logan, E.\ R.\ \textit{et al.}
A study of the physical properties of Li-ion battery electrolytes
containing esters.
\textit{J.\ Electrochem.\ Soc.} \textbf{165}, A21--A30 (2018).

\bibitem{Luo2025}
Luo, Y.\ \textit{et al.}
Solvation configurational entropy governs interfacial kinetics in
metal-ion batteries.
\textit{Joule} \textbf{9}, 210--228 (2025).

\bibitem{Qian2015}
Qian, J.\ \textit{et al.}
High rate and stable cycling of lithium metal anode.
\textit{Nat.\ Commun.} \textbf{6}, 6362 (2015).

\bibitem{Reichardt2011}
Reichardt, C.\ \& Welton, T.
\textit{Solvents and Solvent Effects in Organic Chemistry}, 4th edn.
Wiley-VCH, Weinheim (2011).

\bibitem{Tikekar2016}
Tikekar, M.\ D., Choudhury, S., Tu, Z.\ \& Archer, L.\ A.
Design principles for electrolytes and interfaces for stable
lithium-metal batteries.
\textit{Nat.\ Energy} \textbf{1}, 16114 (2016).

\bibitem{Wan2021}
Wan, G.\ \textit{et al.}
Lithium metal anode operating under lean electrolyte conditions:
a review.
\textit{Adv.\ Mater.} \textbf{33}, 2105713 (2021).

\bibitem{Yu2022}
Yu, Z.\ \textit{et al.}
Rational solvent molecule tuning for high-performance lithium metal
battery electrolytes.
\textit{Nat.\ Energy} \textbf{7}, 94--106 (2022).

\bibitem{Zhang2021}
Zhang, H.\ \textit{et al.}
Fluorinated electrolytes for 5-V lithium-ion battery chemistry.
\textit{Energy Environ.\ Sci.} \textbf{11}, 678--685 (2018).

\end{thebibliography}

\end{document}
